\section{Build e push da primeira image}

\subsection{Docker Hub}

Para compartilhar Docker images, é necessário um local para armazená-las. Nesse contexto entra o registro. Existem vários registros, mas o Docker Hub é o registro padrão para images. O Docker Hub provê tanto um local para armazenar images próprias quanto para encontrar images de outros usuários.

Ao escolher images, o Docker Hub fornece duas categorias de images confiáveis:

\begin{itemize}
    \item \textbf{Docker Official Images (DOI)}: images selecionadas dos softwares mais populares, seguindo boas práticas e atualizadas regularmente.
    \item \textbf{Docker Hardened Images (DHI)}: images desenhadas para reduzir a superfície de ataque --- seguras, simples e praticamente sem vulnerabilidades.
\end{itemize}

Para subir images no Docker Hub, é necessário ter uma conta Docker.

\subsection{Criando um repositório}

Para criar um repositório no Docker Hub, acessa-se o site, seleciona-se a opção de criar repositório e preenche-se as informações necessárias, como o nome do repositório, uma descrição e a visibilidade (pública ou privada).

\subsection{Conceitos fundamentais}

Uma \textbf{image} é um pacote único que contém tudo o que é necessário para rodar um processo. Qualquer máquina que rodar o container usando a image está apta a rodar a aplicação sem precisar instalar nada previamente.

Um \textbf{Dockerfile} é um script que provê o conjunto de instruções de como buildar a image.

\subsection{Comandos de build e push}

Primeiramente, é necessário fazer login no Docker Hub pelo terminal:

\begin{verbatim}
docker login
\end{verbatim}

Para buildar a image:

\begin{verbatim}
docker build -t <DOCKER_USERNAME>/getting-started-todo-app .
\end{verbatim}

Para verificar se a image existe localmente:

\begin{verbatim}
docker image ls
\end{verbatim}

Para subir a image para o Docker Hub:

\begin{verbatim}
docker push <DOCKER_USERNAME>/getting-started-todo-app
\end{verbatim}
